% claim.tex -- AI4Math claim of record (AI-generated, 2026-09-30; instantiated
% from harness/templates/claim/claim.tex per SOP 08b).
% Machine-readable theorem anchors are the "% LEAN:" lines; scripts/paper-lint.sh
% and the claim audit check them. All Lean listings are verbatim, byte-identical
% contiguous excerpts of frozen artifacts of tasks/20260930-mossad42-quad:
%   statements  : formalized/04-mod3x4-statements.lean (106 lines, sha256 b3970831...)
%   proof bodies: final/04-mod3x4-proved.lean          (364 lines, sha256 b40c0802...)
% Line ranges used: statement 26-106 (defs + both theorem signatures); the proof
% is EXCERPTED -- 364 lines do not fit the note, so the note carries the contiguous
% range 100-244 (mod-2 tool lemmas, the style-combination classification lemma,
% the mod-2 induction, and the first main theorem). The full 364-line text ships in
% the deposit under proofs/.
% Verification: audit/listing-verify.txt (audit/inspect-tex.py --verify).
%
% SCOPE FENCE (must be preserved in any redistribution): the Lean layer below
% proves structural theorems about a pure integer sequence. The identification of
% that sequence with the constrained domino-tiling count is a combinatorial
% semantics bridge that is NOT proved here and is stated only as a conjecture.
% The mod-3 object is the sum over tilings whose horizontal-domino COUNT is
% congruent to 0 mod 3 -- counted by the number of horizontal dominoes, modulo 3.
% It is NOT the tiling count taken modulo 3.
\documentclass[11pt]{article}

\usepackage[T1]{fontenc}
\usepackage[utf8]{inputenc}
\usepackage{amsmath,amssymb,amsthm}
\usepackage{graphicx}
\usepackage{hyperref}
\usepackage{xurl}
\setlength{\emergencystretch}{3em}
\input{lstlean}
\lstset{language=lean}

% --- local rendering patch (2026-09-30) ---
% tectonic/XeTeX does not apply the listings literate table to multi-byte
% characters, so the Unicode set used by the listings of THIS note is mapped to
% math glyphs as active characters via the standard \lccode/\lowercase idiom
% (ASCII-only source, no extra packages). Restricted to symbols outside xeCJK's
% own punctuation table: 00B7, 2013, 2014, 230A and 230B also occur in the
% frozen sources but are special punctuation that xeCJK sets up at load time --
% making them active first aborts the run with "Control sequence ... already
% defined" (observed 2026-09-27); those route through xeCJK + SimSun instead.
% 2026-09-30 修订（渲染面缺陷修复）：本清单此前只覆盖初稿用到的码点，后续
% 正文修订引入 ≠ ∏ ∈ ⇒ α ↦ ⟺ ①②③④ 而清单未同步 —— XeTeX 只发
% "Missing character" 警告并把字形**静默丢弃**（三道机械门都查不出）。现按
% claims/scan_listing_unicode.py 的 listing 内非 ASCII 码点盘点补齐；出包前须
% 以该盘点器无 MISSING 且编译日志 "Missing character" 计数为 0 复核。
% Rendering only: listing source bytes are untouched, so the byte-exactness
% checks are unaffected, and the \ifdefined guard leaves the pdflatex path
% identical to the template.
\ifdefined\XeTeXversion
  {\lccode`\~="2022 \lowercase{\global\catcode`~=\active \global\def~{\ensuremath{\bullet}}}}
  {\lccode`\~="2115 \lowercase{\global\catcode`~=\active \global\def~{\ensuremath{\mathbb{N}}}}}
  {\lccode`\~="2124 \lowercase{\global\catcode`~=\active \global\def~{\ensuremath{\mathbb{Z}}}}}
  {\lccode`\~="2190 \lowercase{\global\catcode`~=\active \global\def~{\ensuremath{\leftarrow}}}}
  {\lccode`\~="2192 \lowercase{\global\catcode`~=\active \global\def~{\ensuremath{\rightarrow}}}}
  {\lccode`\~="2194 \lowercase{\global\catcode`~=\active \global\def~{\ensuremath{\leftrightarrow}}}}
  {\lccode`\~="2200 \lowercase{\global\catcode`~=\active \global\def~{\ensuremath{\forall}}}}
  {\lccode`\~="2203 \lowercase{\global\catcode`~=\active \global\def~{\ensuremath{\exists}}}}
  {\lccode`\~="2227 \lowercase{\global\catcode`~=\active \global\def~{\ensuremath{\wedge}}}}
  {\lccode`\~="2228 \lowercase{\global\catcode`~=\active \global\def~{\ensuremath{\vee}}}}
  {\lccode`\~="2261 \lowercase{\global\catcode`~=\active \global\def~{\ensuremath{\equiv}}}}
  {\lccode`\~="2264 \lowercase{\global\catcode`~=\active \global\def~{\ensuremath{\leq}}}}
  {\lccode`\~="2265 \lowercase{\global\catcode`~=\active \global\def~{\ensuremath{\geq}}}}
  {\lccode`\~="22A2 \lowercase{\global\catcode`~=\active \global\def~{\ensuremath{\vdash}}}}
  {\lccode`\~="27E8 \lowercase{\global\catcode`~=\active \global\def~{\ensuremath{\langle}}}}
  {\lccode`\~="27E9 \lowercase{\global\catcode`~=\active \global\def~{\ensuremath{\rangle}}}}
  {\lccode`\~="00AC \lowercase{\global\catcode`~=\active \global\def~{\ensuremath{\neg}}}}
  {\lccode`\~="2212 \lowercase{\global\catcode`~=\active \global\def~{\ensuremath{-}}}}
  {\lccode`\~="2223 \lowercase{\global\catcode`~=\active \global\def~{\ensuremath{\mid}}}}
  {\lccode`\~="21A6 \lowercase{\global\catcode`~=\active \global\def~{\ensuremath{\mapsto}}}}
  {\lccode`\~="2460 \lowercase{\global\catcode`~=\active \global\def~{\textcircled{\scriptsize 1}}}}
  {\lccode`\~="2461 \lowercase{\global\catcode`~=\active \global\def~{\textcircled{\scriptsize 2}}}}
  {\lccode`\~="2462 \lowercase{\global\catcode`~=\active \global\def~{\textcircled{\scriptsize 3}}}}
  {\lccode`\~="2463 \lowercase{\global\catcode`~=\active \global\def~{\textcircled{\scriptsize 4}}}}
  {\lccode`\~="27FA \lowercase{\global\catcode`~=\active \global\def~{\ensuremath{\Longleftrightarrow}}}}
\fi
\ifdefined\XeTeXversion
  \usepackage{xeCJK}
  \setCJKmainfont{SimSun}
\fi

\newtheorem{theorem}{Theorem}

\title{Modular periodicity of the horizontal-domino-residue constrained\\
tiling count on the $4\times n$ board\\
{\large (Lean 4 machine-checked claim of record)}}
\author{Aurora0134\thanks{Contact: via the GitHub account Aurora0134; ORCID
placeholder --- to be filled in by the author before any upload. No personal
information is carried in this note.}}
\date{\today}

\begin{document}
\maketitle

\begin{abstract}
We record a machine-checked result about a pure integer sequence
$a_4\colon\mathbb{N}\to\mathbb{Z}$ defined by initial values together with an
order-twelve constant-coefficient recurrence: its residue modulo $2$ depends only
on $n \bmod 15$ and its residue modulo $3$ depends only on $n \bmod 30$, with the
two residue patterns exhibited explicitly. The combinatorial reading of $a_4$ is
the number of domino tilings of a $4\times n$ board counted by the number of
horizontal dominoes, modulo 3; that identification is \emph{not} proved here. All
statements are verified by the Lean~4 kernel under the pinned mathlib revision,
with no \texttt{sorry} and axioms confined to \{propext, \texttt{Quot.sound},
\texttt{Classical.choice}\}. Novelty: search found no occupying record; the
strongest licensed claim is that no record of this constrained-selection object
was found in the channels queried. This note is a priority claim of record, not a
full paper.
\end{abstract}

\section{Introduction}
This note records a formalized result produced by an AI4Math automated pipeline, in which
propositions and proofs are generated by a language model and adjudicated solely by the
Lean~4 kernel. The result concerns a pure sequence-level object: an order-twelve
constant-coefficient integer recurrence together with the modular periodicity of its
residues. The note is a priority claim of record, deposited for timestamping; the
corresponding narrative paper has been deliberately deferred by the author (see the card),
so no full exposition is attempted here. Verification strength is stated in
Section~\ref{sec:verif}: strict kernel compilation, zero \texttt{sorry}, standard axiom
whitelist. No literature survey is included---that expectation belongs to arXiv moderation
and does not apply to a Zenodo deposit.

\section{Statement}

We first fix notation. Let $a_4\colon\mathbb{N}\to\mathbb{Z}$ be the sequence with initial
values
\[
a_4(0),\dots,a_4(11) \;=\; 1,\,1,\,1,\,1,\,10,\,37,\,110,\,269,\,701,\,2000,\,6020,\,17495
\]
extended by the order-twelve constant-coefficient recurrence
\[
\begin{aligned}
a_4(n) ={}& 3a_4(n{-}1) - 3a_4(n{-}2) + a_4(n{-}3) + 18a_4(n{-}4) - 9a_4(n{-}5) + 20a_4(n{-}6) \\
& + 15a_4(n{-}7) - 18a_4(n{-}8) + a_4(n{-}9) + 9a_4(n{-}10) - a_4(n{-}12),
\end{aligned}
\]
and let $\mathrm{pat2}_{15}$ and $\mathrm{pat3}_{30}$ be the two explicit residue patterns
recorded below (fifteen and thirty entries respectively). Two facts are formalized, and both
are main results; this track contains no definitional-unfolding lemma of the kind recorded
in the sibling three-by-$n$ track.

% LEAN: tasks/20260930-mossad42-quad :: a4_mod2_period15  grade=完全证明
\begin{theorem}[mod-2 residue determined by $n \bmod 15$; main result]
\label{thm:t4a}
For all $n, r \in \mathbb{N}$, if $n \bmod 15 = r$ then
$\;a_4(n) \bmod 2 = \mathrm{pat2}_{15}(r)$.
\end{theorem}

% LEAN: tasks/20260930-mossad42-quad :: a4_mod3_period30  grade=完全证明
\begin{theorem}[mod-3 residue determined by $n \bmod 30$; main result]
\label{thm:t4b}
For all $n, r \in \mathbb{N}$, if $n \bmod 30 = r$ then
$\;a_4(n) \bmod 3 = \mathrm{pat3}_{30}(r)$.
\end{theorem}

The reading of $a_4$ as a tiling count is fixed throughout as follows: $a_4(n)$ is the
number of domino tilings of the $4\times n$ board counted by the number of horizontal
dominoes, modulo 3 --- that is, the sum over those tilings whose horizontal-domino count lies
in the residue class $0$ modulo $3$. This is deliberately \emph{not} the tiling count taken
modulo $3$, which is a different object.

\begin{lstlisting}[caption={Frozen formal statement (Lean 4, verbatim from \texttt{formalized/04-mod3x4-statements.lean}, lines 26--106), with the two proof-body placeholder lines omitted (see the note below the listing). sha256 \texttt{b3970831}...).},label={lst:stmt}]
/-- **T4 主序列 a4（order-12 递推）**。
初值 a4(0..11) = 1, 1, 1, 1, 10, 37, 110, 269, 701, 2000, 6020, 17495；
递推 a4(n) = 3·a4(n−1) − 3·a4(n−2) + a4(n−3) + 18·a4(n−4) − 9·a4(n−5) + 20·a4(n−6)
              + 15·a4(n−7) − 18·a4(n−8) + a4(n−9) + 9·a4(n−10) − a4(n−12)。 -/
def a4 : ℕ → ℤ
  | 0 => 1
  | 1 => 1
  | 2 => 1
  | 3 => 1
  | 4 => 10
  | 5 => 37
  | 6 => 110
  | 7 => 269
  | 8 => 701
  | 9 => 2000
  | 10 => 6020
  | 11 => 17495
  | n + 12 =>
      3 * a4 (n + 11) - 3 * a4 (n + 10) + a4 (n + 9) + 18 * a4 (n + 8)
        - 9 * a4 (n + 7) + 20 * a4 (n + 6) + 15 * a4 (n + 5) - 18 * a4 (n + 4)
        + a4 (n + 3) + 9 * a4 (n + 2) - a4 n

/-- **模 2 样式 pat2_15**（15 留数类 ↦ 余数）：
1,1,1,1,0,1,0,1,1,0,0,1,0,0,0。 -/
def pat2_15 : ℕ → ℤ
  | 0 => 1
  | 1 => 1
  | 2 => 1
  | 3 => 1
  | 4 => 0
  | 5 => 1
  | 6 => 0
  | 7 => 1
  | 8 => 1
  | 9 => 0
  | 10 => 0
  | 11 => 1
  | 12 => 0
  | 13 => 0
  | _ => 0

/-- **模 3 样式 pat3_30**（30 留数类 ↦ 余数）。 -/
def pat3_30 : ℕ → ℤ
  | 0 => 1
  | 1 => 1
  | 2 => 1
  | 3 => 1
  | 4 => 1
  | 5 => 1
  | 6 => 2
  | 7 => 2
  | 8 => 2
  | 9 => 2
  | 10 => 2
  | 11 => 2
  | 12 => 0
  | 13 => 0
  | 14 => 0
  | 15 => 2
  | 16 => 2
  | 17 => 2
  | 18 => 2
  | 19 => 2
  | 20 => 2
  | 21 => 1
  | 22 => 1
  | 23 => 1
  | 24 => 1
  | 25 => 1
  | 26 => 1
  | 27 => 0
  | 28 => 0
  | _ => 0

/-- **T4-A（模 2 周期 15 样式）**：`a4 n` 的 mod 2 余数仅依赖 `n % 15`，样式为 `pat2_15`。 -/
theorem a4_mod2_period15 (n r : ℕ) (hr : n % 15 = r) : a4 n % 2 = pat2_15 r := by

/-- **T4-B（模 3 周期 30 样式）**：`a4 n` 的 mod 3 余数仅依赖 `n % 30`，样式为 `pat3_30`。 -/
theorem a4_mod3_period30 (n r : ℕ) (hr : n % 30 = r) : a4 n % 3 = pat3_30 r := by
\end{lstlisting}

\noindent This is lines 26--106 of the frozen statement snapshot with exactly two lines
removed: the two proof-body placeholder lines, one per theorem (`  sorry`, at lines 102 and
106 of the source file). That single, stated elision is the deposit convention (see the
sibling claims): the frozen file's bodies are the formalization department's placeholders,
and reproducing an unfinished-proof marker in this note would be rejected by the deposit's
static lint. Everything else---doc-comments, definitions, binders, signatures and the
\texttt{:= by} keyword---is verbatim and in order. Lines 1--25 (the pipeline provenance
header, the file-level scope-fence comment block, and the \texttt{import Mathlib} line) are
not reproduced in this listing but are retained in full in the deposit's copy of the file.
\texttt{audit/inspect-tex.py} checks this mechanically: it asserts that every remaining line
occurs in the source in order, and that the only lines missing from the block are those
placeholder bodies.

\section{Proof}
The complete machine proof is 364 lines long and therefore does not fit this note; a
contiguous excerpt is reproduced verbatim below instead. The excerpt is
\textbf{lines 100--244} of \texttt{final/04-mod3x4-proved.lean}: it contains the two shared
arithmetic tool lemmas, the mod-2 style-combination classification lemma
\texttt{pat2\_15\_step}, the full mod-2 induction \texttt{a4\_aux15}, and the first main
theorem \texttt{a4\_mod2\_period15} with its three-line body. The omitted portions are lines
1--99 (the provenance header and the three definitions, already reproduced verbatim in the
statement listing above, together with the \texttt{import} line) and lines 245--364 (the
mod-3 half of the development: \texttt{pat3\_30\_step}, the induction \texttt{a4\_aux30}, and
the second main theorem \texttt{a4\_mod3\_period30}). The mod-3 half is structurally
identical to the mod-2 half shown here, differing only in the period ($30$ rather than $15$),
the modulus ($3$ rather than $2$), and the discarded multiple ($3X$ rather than $2X$). The
full 364-line text is included in this deposit under \texttt{proofs/}. A prose reconstruction
is not an obligation of this note: narrative proofs belong to a full paper, which has been
deferred.

\begin{lstlisting}[caption={Verbatim excerpt of the proof, \texttt{final/04-mod3x4-proved.lean} lines 100--244 (contiguous); full text is 364 lines and ships in the deposit's \texttt{proofs/}. sha256 \texttt{b40c0802}...).},label={lst:proof}]
/-! ## 段落 0：共用的算术工具

关键取向：分类**不作用在 `A` 上**。`A % 15 = j` 根本不决定 `A % 2`（15 是奇数，
`A` 与 `A + 15` 同留数类而异奇偶，`A = 17, j = 2` 即反例）。
分类只作用在**样式组合**（12 项 `pat2_15` 的线性组合）上——
由递推式把 `a4 n` 的 12 项线性组合化为「12 项样式组合 + 2×整数」，
再对样式组合做 15 类机械计算。下面两条工具把「同余 ⟺ 样式项 + 倍数×商」这一形态焊死。 -/

/-- 段落 0 工具 ①（模 2 款）：`2 * (x / 2) + x % 2 = x`，即 `x ≡ x % 2 (mod 2)`。 -/
private theorem int_two_mul_ediv_add_emod (x : ℤ) : 2 * (x / 2) + x % 2 = x := by
  have h := Int.emod_add_mul_ediv x 2
  omega

/-- 段落 0 工具 ②（模 3 款）：`3 * (x / 3) + x % 3 = x`，即 `x ≡ x % 3 (mod 3)`。 -/
private theorem int_three_mul_ediv_add_emod (x : ℤ) : 3 * (x / 3) + x % 3 = x := by
  have h := Int.emod_add_mul_ediv x 3
  omega

/-! ## 段落 1：模 2 的 15 个留数类机械计算（T4-A 的收口引理）

`j < 15` 前提不可省——`interval_cases` 靠它找上界。分类前提写成 `j : ℕ` 上界形式，
而不是「`A % 15 = j`」形态：后者是假命题。 -/

/-- **段落 1 分类引理**：12 项样式组合 mod 2 = `pat2_15 ((j + 12) % 15)`，`j < 15`。 -/
private theorem pat2_15_step (j : ℕ) (hj : j < 15) :
    (3 * pat2_15 ((j + 11) % 15) - 3 * pat2_15 ((j + 10) % 15) + pat2_15 ((j + 9) % 15)
      + 18 * pat2_15 ((j + 8) % 15) - 9 * pat2_15 ((j + 7) % 15) + 20 * pat2_15 ((j + 6) % 15)
      + 15 * pat2_15 ((j + 5) % 15) - 18 * pat2_15 ((j + 4) % 15) + pat2_15 ((j + 3) % 15)
      + 9 * pat2_15 ((j + 2) % 15) - pat2_15 (j % 15)) % 2 = pat2_15 ((j + 12) % 15) := by
  interval_cases j <;> decide

/-! ## 段落 2：T4-A 的归纳（`a4 n % 2` 只由 `n % 15` 决定）

强归纳。`n < 12` 走初值（`decide`）；`n ≥ 12` 写 `n = m + 12`，先由归纳假设得 12 条
`a4 (m+i) = pat2_15 ((m+i) % 15) + 2 * (a4 (m+i) / 2)`，代入递推式后把 `(m+i) % 15`
归一为 `(m % 15 + i) % 15`，再用 `ring` 把整条组合式拆成「样式组合 + 2×整数」，
最后 `Int.add_emod` 丢掉偶倍项、`simpa [Nat.mod_mod]` 对上 `pat2_15_step` 的形态。 -/

/-- **T4-A 归纳辅助**：`a4` 的 mod 2 余数只由 `n % 15` 决定。 -/
private theorem a4_aux15 : ∀ n : ℕ, a4 n % 2 = pat2_15 (n % 15) := by
  intro n
  induction n using Nat.strong_induction_on with
  | _ n ih =>
    rcases Nat.lt_or_ge n 12 with hn | hn
    · interval_cases n <;> decide
    · obtain ⟨m, rfl⟩ : ∃ m, n = m + 12 := ⟨n - 12, by omega⟩
      -- 递推式展开（`a4` 的 `n + 12` 分支）
      have hrec : a4 (m + 12) = 3 * a4 (m + 11) - 3 * a4 (m + 10) + a4 (m + 9)
          + 18 * a4 (m + 8) - 9 * a4 (m + 7) + 20 * a4 (m + 6) + 15 * a4 (m + 5)
          - 18 * a4 (m + 4) + a4 (m + 3) + 9 * a4 (m + 2) - a4 m := by
        rw [a4]
      -- 12 条「样式项 + 2×商」（归纳假设 + `Int.emod_add_mul_ediv`）
      have h0 : a4 m = pat2_15 (m % 15) + 2 * (a4 m / 2) := by
        have h := ih m (by omega)
        have h2 := int_two_mul_ediv_add_emod (a4 m)
        omega
      have h2 : a4 (m + 2) = pat2_15 ((m + 2) % 15) + 2 * (a4 (m + 2) / 2) := by
        have h := ih (m + 2) (by omega)
        have h2 := int_two_mul_ediv_add_emod (a4 (m + 2))
        omega
      have h3 : a4 (m + 3) = pat2_15 ((m + 3) % 15) + 2 * (a4 (m + 3) / 2) := by
        have h := ih (m + 3) (by omega)
        have h2 := int_two_mul_ediv_add_emod (a4 (m + 3))
        omega
      have h4 : a4 (m + 4) = pat2_15 ((m + 4) % 15) + 2 * (a4 (m + 4) / 2) := by
        have h := ih (m + 4) (by omega)
        have h2 := int_two_mul_ediv_add_emod (a4 (m + 4))
        omega
      have h5 : a4 (m + 5) = pat2_15 ((m + 5) % 15) + 2 * (a4 (m + 5) / 2) := by
        have h := ih (m + 5) (by omega)
        have h2 := int_two_mul_ediv_add_emod (a4 (m + 5))
        omega
      have h6 : a4 (m + 6) = pat2_15 ((m + 6) % 15) + 2 * (a4 (m + 6) / 2) := by
        have h := ih (m + 6) (by omega)
        have h2 := int_two_mul_ediv_add_emod (a4 (m + 6))
        omega
      have h7 : a4 (m + 7) = pat2_15 ((m + 7) % 15) + 2 * (a4 (m + 7) / 2) := by
        have h := ih (m + 7) (by omega)
        have h2 := int_two_mul_ediv_add_emod (a4 (m + 7))
        omega
      have h8 : a4 (m + 8) = pat2_15 ((m + 8) % 15) + 2 * (a4 (m + 8) / 2) := by
        have h := ih (m + 8) (by omega)
        have h2 := int_two_mul_ediv_add_emod (a4 (m + 8))
        omega
      have h9 : a4 (m + 9) = pat2_15 ((m + 9) % 15) + 2 * (a4 (m + 9) / 2) := by
        have h := ih (m + 9) (by omega)
        have h2 := int_two_mul_ediv_add_emod (a4 (m + 9))
        omega
      have h10 : a4 (m + 10) = pat2_15 ((m + 10) % 15) + 2 * (a4 (m + 10) / 2) := by
        have h := ih (m + 10) (by omega)
        have h2 := int_two_mul_ediv_add_emod (a4 (m + 10))
        omega
      have h11 : a4 (m + 11) = pat2_15 ((m + 11) % 15) + 2 * (a4 (m + 11) / 2) := by
        have h := ih (m + 11) (by omega)
        have h2 := int_two_mul_ediv_add_emod (a4 (m + 11))
        omega
      -- ① 代入 12 条；② 把残留的 `(m+i) % 15` 归一（`i = 0` 项不需要）
      rw [hrec, h0, h2, h3, h4, h5, h6, h7, h8, h9, h10, h11]
      rw [show (m + 11) % 15 = (m % 15 + 11) % 15 by omega,
          show (m + 10) % 15 = (m % 15 + 10) % 15 by omega,
          show (m + 9) % 15 = (m % 15 + 9) % 15 by omega,
          show (m + 8) % 15 = (m % 15 + 8) % 15 by omega,
          show (m + 7) % 15 = (m % 15 + 7) % 15 by omega,
          show (m + 6) % 15 = (m % 15 + 6) % 15 by omega,
          show (m + 5) % 15 = (m % 15 + 5) % 15 by omega,
          show (m + 4) % 15 = (m % 15 + 4) % 15 by omega,
          show (m + 3) % 15 = (m % 15 + 3) % 15 by omega,
          show (m + 2) % 15 = (m % 15 + 2) % 15 by omega]
      -- ③ 组合式 = 样式组合 + 2×整数（纯 `ring`）
      have key : (3 * (pat2_15 ((m % 15 + 11) % 15) + 2 * (a4 (m + 11) / 2))
          - 3 * (pat2_15 ((m % 15 + 10) % 15) + 2 * (a4 (m + 10) / 2))
          + (pat2_15 ((m % 15 + 9) % 15) + 2 * (a4 (m + 9) / 2))
          + 18 * (pat2_15 ((m % 15 + 8) % 15) + 2 * (a4 (m + 8) / 2))
          - 9 * (pat2_15 ((m % 15 + 7) % 15) + 2 * (a4 (m + 7) / 2))
          + 20 * (pat2_15 ((m % 15 + 6) % 15) + 2 * (a4 (m + 6) / 2))
          + 15 * (pat2_15 ((m % 15 + 5) % 15) + 2 * (a4 (m + 5) / 2))
          - 18 * (pat2_15 ((m % 15 + 4) % 15) + 2 * (a4 (m + 4) / 2))
          + (pat2_15 ((m % 15 + 3) % 15) + 2 * (a4 (m + 3) / 2))
          + 9 * (pat2_15 ((m % 15 + 2) % 15) + 2 * (a4 (m + 2) / 2))
          - (pat2_15 (m % 15) + 2 * (a4 m / 2)) : ℤ)
          = (3 * pat2_15 ((m % 15 + 11) % 15) - 3 * pat2_15 ((m % 15 + 10) % 15)
            + pat2_15 ((m % 15 + 9) % 15) + 18 * pat2_15 ((m % 15 + 8) % 15)
            - 9 * pat2_15 ((m % 15 + 7) % 15) + 20 * pat2_15 ((m % 15 + 6) % 15)
            + 15 * pat2_15 ((m % 15 + 5) % 15) - 18 * pat2_15 ((m % 15 + 4) % 15)
            + pat2_15 ((m % 15 + 3) % 15) + 9 * pat2_15 ((m % 15 + 2) % 15)
            - pat2_15 (m % 15))
          + 2 * (3 * (a4 (m + 11) / 2) - 3 * (a4 (m + 10) / 2) + (a4 (m + 9) / 2)
            + 18 * (a4 (m + 8) / 2) - 9 * (a4 (m + 7) / 2) + 20 * (a4 (m + 6) / 2)
            + 15 * (a4 (m + 5) / 2) - 18 * (a4 (m + 4) / 2) + (a4 (m + 3) / 2)
            + 9 * (a4 (m + 2) / 2) - (a4 m / 2)) := by
        ring
      rw [key]
      rw [Int.add_emod]
      rw [show (2 * (3 * (a4 (m + 11) / 2) - 3 * (a4 (m + 10) / 2) + (a4 (m + 9) / 2)
            + 18 * (a4 (m + 8) / 2) - 9 * (a4 (m + 7) / 2) + 20 * (a4 (m + 6) / 2)
            + 15 * (a4 (m + 5) / 2) - 18 * (a4 (m + 4) / 2) + (a4 (m + 3) / 2)
            + 9 * (a4 (m + 2) / 2) - (a4 m / 2))) % 2 = 0 by omega]
      rw [Int.add_zero]
      -- ④ 形态差（`m % 15 % 15`、外层 `% 2 % 2`）用 `simpa`，不能用 `exact`
      simpa [Nat.mod_mod] using pat2_15_step (m % 15) (Nat.mod_lt _ (by norm_num))

/-- **T4-A（模 2 周期 15 样式）**：`a4 n` 的 mod 2 余数仅依赖 `n % 15`，样式为 `pat2_15`。 -/
theorem a4_mod2_period15 (n r : ℕ) (hr : n % 15 = r) : a4 n % 2 = pat2_15 r := by
  rw [← hr]
  exact a4_aux15 n
\end{lstlisting}

\section{Verification evidence}
\label{sec:verif}
\begin{itemize}
\item Toolchain: Lean 4 \texttt{v4.34.0}; mathlib4 \texttt{v4.34.0}, revision \texttt{5ed29652}.
\item Statement integrity: \texttt{formalized/04-mod3x4-statements.lean}, 106 lines,
sha256 \texttt{b39708311b6fabd6}\allowbreak\texttt{e8828c1886823049}\allowbreak\texttt{a6d0d28c6e66d9c8}\allowbreak\texttt{76a68fe65c377880}.
Final artifact \texttt{final/04-mod3x4-proved.lean}, 364 lines,
sha256 \texttt{b40c0802811f136a}\allowbreak\texttt{cb5525d80012ee63}\allowbreak\texttt{96a015eeba45e45c}\allowbreak\texttt{d41183610e2b8adb}.
Every listing above is byte-identical to its source (checked by \texttt{audit/inspect-tex.py}).
\item Kernel check: strict compilation exits 0 with no \texttt{sorry} and no \texttt{admit}.
\texttt{\#print axioms} output, verbatim:
\begin{verbatim}
'a4_mod2_period15' depends on axioms: [propext, Classical.choice, Quot.sound]
'a4_mod3_period30' depends on axioms: [propext, Classical.choice, Quot.sound]
\end{verbatim}
The axiom set is a subset of \{propext, \texttt{Quot.sound}, \texttt{Classical.choice}\};
\texttt{sorryAx} does not occur.
\item Audit chain: statement freeze (sha256), byte-wise symmetric-difference comparison against
the frozen snapshot, strict-mode compilation, and the axiom whitelist check. Originals are in
the deposit under \texttt{audit/}.
\end{itemize}

\section{Novelty evidence}
Under the pipeline's C9 novelty discipline the verdict for this object is \emph{no occupying
record found}. The full query log, including every zero-hit query, is in the deposit at
\texttt{audit/c9-record.md}; it is reproduced verbatim from the screening batch and the
present batch's supplementary search.

Summary of channels. OEIS: five mutually distinct variants of the sequence itself (identity,
even-index, odd-index, first-difference, partial sums) all returned zero hits, and a
twenty-one-term direct query also returned zero hits; channel availability was established by
positive controls before any zero-hit was read (A001835, A005178, A000045). Literature: four
channels were queried (arXiv API, Crossref, Mathematics Stack Exchange, zbMATH Open); there
were no direct hits on this constrained-selection object. Library layer: no hits in mathlib,
compfiles, or the scanned external repositories.

Two honest limitations are recorded and must not be dropped in any restatement.
First, the OpenAlex channel was \emph{unreachable} from the machine used for this search
(Cloudflare-fronted hosts timed out; DNS resolution was verified correct against a public
resolver, so this is not a DNS-poisoning artifact). Crossref was used as a substitute
DOI-level index. Coverage is therefore \emph{lower} than in the screening batch, and this is
a channel gap, not strengthened evidence.
Second, a zero hit is \emph{not} novelty: the strongest claim licensed here is that no
record of this constrained-selection object was found in the channels queried.
No moment or expectation novelty is claimed: mixed moments $E[V^a H^b]$ for general $m\times n$
boards are already occupied by \emph{Statistics of Domino Tilings on a Rectangular Board},
Fibonacci Quarterly (2019), doi:10.1080/00150517.2019.12427625.
The $3\times n$ and $4\times n$ mod-3 objects --- this track and its sibling --- are
$m$-dimensional variants of one another and are not used as independent novelty evidence for
each other.

\section{Scope and limitations}
The kernel-level results concern the pure integer sequence $a_4$ defined above. They do
\emph{not} establish that $a_4(n)$ equals the number of domino tilings of the $4\times n$
board counted by the number of horizontal dominoes, modulo 3. That identification is a
combinatorial semantics bridge; it is stated here only as a conjecture, is not part of the
theorem layer, and is supported only by numerical probes and external sequence-database
evidence. Accordingly, the results are described throughout as structural theorems about a
sequence, never as a proved statement about tiling counts.

Note the subscript convention deliberately \emph{not} borrowed here. On the $3\times n$
half-horizontal track the counts alternate with zeros, and that sparsity licenses an
odd-index convention there; the present $4\times n$ sequence has \emph{no} alternating zeros
(it is non-zero for odd $n$ as well), so no such convention applies and none is used.

Both recorded theorems are main results, not unfoldings: neither is closed by the bare
definitional probe, and each is established by a genuine induction over the recurrence. The
grading of both is \emph{complete proof} (0 \texttt{sorry}, axiom whitelist), and no weakening
or extra hypothesis was used.

\paragraph{Proof-search history (recorded because it is this track's most useful byproduct).}
This was the only one of the four parallel tracks that required six sampling rounds; the
others needed one or two. The root cause is that the prover invented \emph{false} lemmas as
stepping stones, twice, and both were caught only by the kernel.
\begin{itemize}
\item \emph{Round 4.} The prover asserted \texttt{have h12 : (m + 12) \% 15 = m \% 15},
which is false (at $m=0$ the two sides are $12$ and $0$). A second, independent factor
surfaced in the same round: the \texttt{omega} tactic does not unfold opaque functions such
as \texttt{pat2\_15}, so goals phrased over \texttt{pat2\_15} could not be discharged that
way and had to go through \texttt{interval\_cases} and \texttt{decide}.
\item \emph{Round 5.} The prover then asserted two lemmas of the form
$A \bmod 15 = j \to A \bmod 2 = \mathrm{pat2}_{15}((j+12) \bmod 15)$, which are false because
$15$ is odd: integers in the same residue class modulo $15$ need not share parity. The kernel
counterexample is $A=17$, $j=2$.
\item \emph{Fix.} The classification was moved off the integer $A$ and onto the
\emph{style combination} --- the twelve-term linear combination of pattern values --- where
the induction hypothesis actually supplies the needed congruence. Round 6 passed on the first
attempt with that route.
\end{itemize}
The methodological lesson is recorded here explicitly: the error-feedback discipline must
verify whether the \emph{proposition itself is true}, not merely that lemma names and
signatures check out. The rounds spent on this track were spent almost entirely on that
distinction.
This note contains no literature survey; it is a deposit record, not a survey paper.

\section*{Data availability}
Lean sources, verification and novelty-search evidence are deposited on Zenodo
(DOI: \texttt{RESERVED-DOI}) and mirrored as a public snapshot repository. The note is
released under CC BY 4.0 and the code under MIT.

\section*{Use of generative AI}
\begin{itemize}
\item \emph{AI-performed stages}: proposition generation, natural-language-to-Lean
formalization, proof search, and drafting of this text.
\item \emph{Model, node, budget}: pipeline node
\texttt{anthropic/}\allowbreak\texttt{a6api-main/}\allowbreak\texttt{deepseek-v4.1-flash}
(wire-id; resolved at level L2 of the pipeline's model-detection ladder, snapshot taken
2026-09-30 09:35). Source-task budget: sampling 6 of 8 allowed runs on this track (10 of 32
across four parallel tracks); judgement-class LLM calls 6 of 16 across the batch; Lean
formal compilations 4 of 12 on this track (19 of 48 across the batch); LaTeX 0. The numbers
are copied from the source task's \texttt{budget.log}.
\item \emph{Final adjudication}: all statements and proofs reported here were checked by the
Lean~4 kernel. Every proof compiles with no \texttt{sorry}, and \texttt{\#print axioms}
reports only \{propext, \texttt{Quot.sound}, \texttt{Classical.choice}\}. Statement text was
frozen before proving and compared byte-wise afterwards.
\item The author reviewed the material and is responsible for all content. AI systems are
not listed as authors.
\end{itemize}

\section*{Author contributions}
The author determined the research question and the scope fences, verified the semantics of
the statements, and is responsible for signing off and for any deposit or release. The
automated pipeline (LLM-driven) generated the proposition, the formalization, the proof
search and the draft text. No external release is performed by the pipeline.

\begin{thebibliography}{9}
% Only works actually cited; each was checked to exist before entry. Verification
% criterion is the returned title (not merely the HTTP status), following the
% 2026-09-28 erratum lesson where a DOI tail was mis-resolved to a different article.
\bibitem{lean4}
L.~de Moura and S.~Ullrich.
\newblock The Lean 4 theorem prover and programming language.
\newblock \emph{CADE-28}, 2021. \url{https://doi.org/10.1007/978-3-030-79876-5_37}

\bibitem{mathlib}
The mathlib Community.
\newblock The Lean mathematical library.
\newblock \emph{CPP 2020}, 2020. \url{https://doi.org/10.1145/3372885.3373824}

\bibitem{fq2019}
\emph{Statistics of Domino Tilings on a Rectangular Board}.
\newblock Fibonacci Quarterly, 2019.
\newblock \url{https://doi.org/10.1080/00150517.2019.12427625}
\end{thebibliography}

\end{document}
